\pdfoutput=1
\pdfminorversion=7
\PassOptionsToPackage{table}{xcolor}
\documentclass{article} % For LaTeX2e
\usepackage{iclr2027_conference,times}

%%%%% NEW MATH DEFINITIONS %%%%%

\usepackage{amsmath,amsfonts,bm}

% Mark sections of captions for referring to divisions of figures
\newcommand{\figleft}{{\em (Left)}}
\newcommand{\figcenter}{{\em (Center)}}
\newcommand{\figright}{{\em (Right)}}
\newcommand{\figtop}{{\em (Top)}}
\newcommand{\figbottom}{{\em (Bottom)}}
\newcommand{\captiona}{{\em (a)}}
\newcommand{\captionb}{{\em (b)}}
\newcommand{\captionc}{{\em (c)}}
\newcommand{\captiond}{{\em (d)}}

% Highlight a newly defined term
\newcommand{\newterm}[1]{{\bf #1}}


% Figure reference, lower-case.
\def\figref#1{figure~\ref{#1}}
% Figure reference, capital. For start of sentence
\def\Figref#1{Figure~\ref{#1}}
\def\twofigref#1#2{figures \ref{#1} and \ref{#2}}
\def\quadfigref#1#2#3#4{figures \ref{#1}, \ref{#2}, \ref{#3} and \ref{#4}}
% Section reference, lower-case.
\def\secref#1{section~\ref{#1}}
% Section reference, capital.
\def\Secref#1{Section~\ref{#1}}
% Reference to two sections.
\def\twosecrefs#1#2{sections \ref{#1} and \ref{#2}}
% Reference to three sections.
\def\secrefs#1#2#3{sections \ref{#1}, \ref{#2} and \ref{#3}}
% Reference to an equation, lower-case.
\def\eqref#1{equation~\ref{#1}}
% Reference to an equation, upper case
\def\Eqref#1{Equation~\ref{#1}}
% A raw reference to an equation---avoid using if possible
\def\plaineqref#1{\ref{#1}}
% Reference to a chapter, lower-case.
\def\chapref#1{chapter~\ref{#1}}
% Reference to an equation, upper case.
\def\Chapref#1{Chapter~\ref{#1}}
% Reference to a range of chapters
\def\rangechapref#1#2{chapters\ref{#1}--\ref{#2}}
% Reference to an algorithm, lower-case.
\def\algref#1{algorithm~\ref{#1}}
% Reference to an algorithm, upper case.
\def\Algref#1{Algorithm~\ref{#1}}
\def\twoalgref#1#2{algorithms \ref{#1} and \ref{#2}}
\def\Twoalgref#1#2{Algorithms \ref{#1} and \ref{#2}}
% Reference to a part, lower case
\def\partref#1{part~\ref{#1}}
% Reference to a part, upper case
\def\Partref#1{Part~\ref{#1}}
\def\twopartref#1#2{parts \ref{#1} and \ref{#2}}

\def\ceil#1{\lceil #1 \rceil}
\def\floor#1{\lfloor #1 \rfloor}
\def\1{\bm{1}}
\newcommand{\train}{\mathcal{D}}
\newcommand{\valid}{\mathcal{D_{\mathrm{valid}}}}
\newcommand{\test}{\mathcal{D_{\mathrm{test}}}}

\def\eps{{\epsilon}}


% Random variables
\def\reta{{\textnormal{$\eta$}}}
\def\ra{{\textnormal{a}}}
\def\rb{{\textnormal{b}}}
\def\rc{{\textnormal{c}}}
\def\rd{{\textnormal{d}}}
\def\re{{\textnormal{e}}}
\def\rf{{\textnormal{f}}}
\def\rg{{\textnormal{g}}}
\def\rh{{\textnormal{h}}}
\def\ri{{\textnormal{i}}}
\def\rj{{\textnormal{j}}}
\def\rk{{\textnormal{k}}}
\def\rl{{\textnormal{l}}}
% rm is already a command, just don't name any random variables m
\def\rn{{\textnormal{n}}}
\def\ro{{\textnormal{o}}}
\def\rp{{\textnormal{p}}}
\def\rq{{\textnormal{q}}}
\def\rr{{\textnormal{r}}}
\def\rs{{\textnormal{s}}}
\def\rt{{\textnormal{t}}}
\def\ru{{\textnormal{u}}}
\def\rv{{\textnormal{v}}}
\def\rw{{\textnormal{w}}}
\def\rx{{\textnormal{x}}}
\def\ry{{\textnormal{y}}}
\def\rz{{\textnormal{z}}}

% Random vectors
\def\rvepsilon{{\mathbf{\epsilon}}}
\def\rvtheta{{\mathbf{\theta}}}
\def\rva{{\mathbf{a}}}
\def\rvb{{\mathbf{b}}}
\def\rvc{{\mathbf{c}}}
\def\rvd{{\mathbf{d}}}
\def\rve{{\mathbf{e}}}
\def\rvf{{\mathbf{f}}}
\def\rvg{{\mathbf{g}}}
\def\rvh{{\mathbf{h}}}
\def\rvu{{\mathbf{i}}}
\def\rvj{{\mathbf{j}}}
\def\rvk{{\mathbf{k}}}
\def\rvl{{\mathbf{l}}}
\def\rvm{{\mathbf{m}}}
\def\rvn{{\mathbf{n}}}
\def\rvo{{\mathbf{o}}}
\def\rvp{{\mathbf{p}}}
\def\rvq{{\mathbf{q}}}
\def\rvr{{\mathbf{r}}}
\def\rvs{{\mathbf{s}}}
\def\rvt{{\mathbf{t}}}
\def\rvu{{\mathbf{u}}}
\def\rvv{{\mathbf{v}}}
\def\rvw{{\mathbf{w}}}
\def\rvx{{\mathbf{x}}}
\def\rvy{{\mathbf{y}}}
\def\rvz{{\mathbf{z}}}

% Elements of random vectors
\def\erva{{\textnormal{a}}}
\def\ervb{{\textnormal{b}}}
\def\ervc{{\textnormal{c}}}
\def\ervd{{\textnormal{d}}}
\def\erve{{\textnormal{e}}}
\def\ervf{{\textnormal{f}}}
\def\ervg{{\textnormal{g}}}
\def\ervh{{\textnormal{h}}}
\def\ervi{{\textnormal{i}}}
\def\ervj{{\textnormal{j}}}
\def\ervk{{\textnormal{k}}}
\def\ervl{{\textnormal{l}}}
\def\ervm{{\textnormal{m}}}
\def\ervn{{\textnormal{n}}}
\def\ervo{{\textnormal{o}}}
\def\ervp{{\textnormal{p}}}
\def\ervq{{\textnormal{q}}}
\def\ervr{{\textnormal{r}}}
\def\ervs{{\textnormal{s}}}
\def\ervt{{\textnormal{t}}}
\def\ervu{{\textnormal{u}}}
\def\ervv{{\textnormal{v}}}
\def\ervw{{\textnormal{w}}}
\def\ervx{{\textnormal{x}}}
\def\ervy{{\textnormal{y}}}
\def\ervz{{\textnormal{z}}}

% Random matrices
\def\rmA{{\mathbf{A}}}
\def\rmB{{\mathbf{B}}}
\def\rmC{{\mathbf{C}}}
\def\rmD{{\mathbf{D}}}
\def\rmE{{\mathbf{E}}}
\def\rmF{{\mathbf{F}}}
\def\rmG{{\mathbf{G}}}
\def\rmH{{\mathbf{H}}}
\def\rmI{{\mathbf{I}}}
\def\rmJ{{\mathbf{J}}}
\def\rmK{{\mathbf{K}}}
\def\rmL{{\mathbf{L}}}
\def\rmM{{\mathbf{M}}}
\def\rmN{{\mathbf{N}}}
\def\rmO{{\mathbf{O}}}
\def\rmP{{\mathbf{P}}}
\def\rmQ{{\mathbf{Q}}}
\def\rmR{{\mathbf{R}}}
\def\rmS{{\mathbf{S}}}
\def\rmT{{\mathbf{T}}}
\def\rmU{{\mathbf{U}}}
\def\rmV{{\mathbf{V}}}
\def\rmW{{\mathbf{W}}}
\def\rmX{{\mathbf{X}}}
\def\rmY{{\mathbf{Y}}}
\def\rmZ{{\mathbf{Z}}}

% Elements of random matrices
\def\ermA{{\textnormal{A}}}
\def\ermB{{\textnormal{B}}}
\def\ermC{{\textnormal{C}}}
\def\ermD{{\textnormal{D}}}
\def\ermE{{\textnormal{E}}}
\def\ermF{{\textnormal{F}}}
\def\ermG{{\textnormal{G}}}
\def\ermH{{\textnormal{H}}}
\def\ermI{{\textnormal{I}}}
\def\ermJ{{\textnormal{J}}}
\def\ermK{{\textnormal{K}}}
\def\ermL{{\textnormal{L}}}
\def\ermM{{\textnormal{M}}}
\def\ermN{{\textnormal{N}}}
\def\ermO{{\textnormal{O}}}
\def\ermP{{\textnormal{P}}}
\def\ermQ{{\textnormal{Q}}}
\def\ermR{{\textnormal{R}}}
\def\ermS{{\textnormal{S}}}
\def\ermT{{\textnormal{T}}}
\def\ermU{{\textnormal{U}}}
\def\ermV{{\textnormal{V}}}
\def\ermW{{\textnormal{W}}}
\def\ermX{{\textnormal{X}}}
\def\ermY{{\textnormal{Y}}}
\def\ermZ{{\textnormal{Z}}}

% Vectors
\def\vzero{{\bm{0}}}
\def\vone{{\bm{1}}}
\def\vmu{{\bm{\mu}}}
\def\vtheta{{\bm{\theta}}}
\def\va{{\bm{a}}}
\def\vb{{\bm{b}}}
\def\vc{{\bm{c}}}
\def\vd{{\bm{d}}}
\def\ve{{\bm{e}}}
\def\vf{{\bm{f}}}
\def\vg{{\bm{g}}}
\def\vh{{\bm{h}}}
\def\vi{{\bm{i}}}
\def\vj{{\bm{j}}}
\def\vk{{\bm{k}}}
\def\vl{{\bm{l}}}
\def\vm{{\bm{m}}}
\def\vn{{\bm{n}}}
\def\vo{{\bm{o}}}
\def\vp{{\bm{p}}}
\def\vq{{\bm{q}}}
\def\vr{{\bm{r}}}
\def\vs{{\bm{s}}}
\def\vt{{\bm{t}}}
\def\vu{{\bm{u}}}
\def\vv{{\bm{v}}}
\def\vw{{\bm{w}}}
\def\vx{{\bm{x}}}
\def\vy{{\bm{y}}}
\def\vz{{\bm{z}}}

% Elements of vectors
\def\evalpha{{\alpha}}
\def\evbeta{{\beta}}
\def\evepsilon{{\epsilon}}
\def\evlambda{{\lambda}}
\def\evomega{{\omega}}
\def\evmu{{\mu}}
\def\evpsi{{\psi}}
\def\evsigma{{\sigma}}
\def\evtheta{{\theta}}
\def\eva{{a}}
\def\evb{{b}}
\def\evc{{c}}
\def\evd{{d}}
\def\eve{{e}}
\def\evf{{f}}
\def\evg{{g}}
\def\evh{{h}}
\def\evi{{i}}
\def\evj{{j}}
\def\evk{{k}}
\def\evl{{l}}
\def\evm{{m}}
\def\evn{{n}}
\def\evo{{o}}
\def\evp{{p}}
\def\evq{{q}}
\def\evr{{r}}
\def\evs{{s}}
\def\evt{{t}}
\def\evu{{u}}
\def\evv{{v}}
\def\evw{{w}}
\def\evx{{x}}
\def\evy{{y}}
\def\evz{{z}}

% Matrix
\def\mA{{\bm{A}}}
\def\mB{{\bm{B}}}
\def\mC{{\bm{C}}}
\def\mD{{\bm{D}}}
\def\mE{{\bm{E}}}
\def\mF{{\bm{F}}}
\def\mG{{\bm{G}}}
\def\mH{{\bm{H}}}
\def\mI{{\bm{I}}}
\def\mJ{{\bm{J}}}
\def\mK{{\bm{K}}}
\def\mL{{\bm{L}}}
\def\mM{{\bm{M}}}
\def\mN{{\bm{N}}}
\def\mO{{\bm{O}}}
\def\mP{{\bm{P}}}
\def\mQ{{\bm{Q}}}
\def\mR{{\bm{R}}}
\def\mS{{\bm{S}}}
\def\mT{{\bm{T}}}
\def\mU{{\bm{U}}}
\def\mV{{\bm{V}}}
\def\mW{{\bm{W}}}
\def\mX{{\bm{X}}}
\def\mY{{\bm{Y}}}
\def\mZ{{\bm{Z}}}
\def\mBeta{{\bm{\beta}}}
\def\mPhi{{\bm{\Phi}}}
\def\mLambda{{\bm{\Lambda}}}
\def\mSigma{{\bm{\Sigma}}}

% Tensor
\DeclareMathAlphabet{\mathsfit}{\encodingdefault}{\sfdefault}{m}{sl}
\SetMathAlphabet{\mathsfit}{bold}{\encodingdefault}{\sfdefault}{bx}{n}
\newcommand{\tens}[1]{\bm{\mathsfit{#1}}}
\def\tA{{\tens{A}}}
\def\tB{{\tens{B}}}
\def\tC{{\tens{C}}}
\def\tD{{\tens{D}}}
\def\tE{{\tens{E}}}
\def\tF{{\tens{F}}}
\def\tG{{\tens{G}}}
\def\tH{{\tens{H}}}
\def\tI{{\tens{I}}}
\def\tJ{{\tens{J}}}
\def\tK{{\tens{K}}}
\def\tL{{\tens{L}}}
\def\tM{{\tens{M}}}
\def\tN{{\tens{N}}}
\def\tO{{\tens{O}}}
\def\tP{{\tens{P}}}
\def\tQ{{\tens{Q}}}
\def\tR{{\tens{R}}}
\def\tS{{\tens{S}}}
\def\tT{{\tens{T}}}
\def\tU{{\tens{U}}}
\def\tV{{\tens{V}}}
\def\tW{{\tens{W}}}
\def\tX{{\tens{X}}}
\def\tY{{\tens{Y}}}
\def\tZ{{\tens{Z}}}


% Graph
\def\gA{{\mathcal{A}}}
\def\gB{{\mathcal{B}}}
\def\gC{{\mathcal{C}}}
\def\gD{{\mathcal{D}}}
\def\gE{{\mathcal{E}}}
\def\gF{{\mathcal{F}}}
\def\gG{{\mathcal{G}}}
\def\gH{{\mathcal{H}}}
\def\gI{{\mathcal{I}}}
\def\gJ{{\mathcal{J}}}
\def\gK{{\mathcal{K}}}
\def\gL{{\mathcal{L}}}
\def\gM{{\mathcal{M}}}
\def\gN{{\mathcal{N}}}
\def\gO{{\mathcal{O}}}
\def\gP{{\mathcal{P}}}
\def\gQ{{\mathcal{Q}}}
\def\gR{{\mathcal{R}}}
\def\gS{{\mathcal{S}}}
\def\gT{{\mathcal{T}}}
\def\gU{{\mathcal{U}}}
\def\gV{{\mathcal{V}}}
\def\gW{{\mathcal{W}}}
\def\gX{{\mathcal{X}}}
\def\gY{{\mathcal{Y}}}
\def\gZ{{\mathcal{Z}}}

% Sets
\def\sA{{\mathbb{A}}}
\def\sB{{\mathbb{B}}}
\def\sC{{\mathbb{C}}}
\def\sD{{\mathbb{D}}}
% Don't use a set called E, because this would be the same as our symbol
% for expectation.
\def\sF{{\mathbb{F}}}
\def\sG{{\mathbb{G}}}
\def\sH{{\mathbb{H}}}
\def\sI{{\mathbb{I}}}
\def\sJ{{\mathbb{J}}}
\def\sK{{\mathbb{K}}}
\def\sL{{\mathbb{L}}}
\def\sM{{\mathbb{M}}}
\def\sN{{\mathbb{N}}}
\def\sO{{\mathbb{O}}}
\def\sP{{\mathbb{P}}}
\def\sQ{{\mathbb{Q}}}
\def\sR{{\mathbb{R}}}
\def\sS{{\mathbb{S}}}
\def\sT{{\mathbb{T}}}
\def\sU{{\mathbb{U}}}
\def\sV{{\mathbb{V}}}
\def\sW{{\mathbb{W}}}
\def\sX{{\mathbb{X}}}
\def\sY{{\mathbb{Y}}}
\def\sZ{{\mathbb{Z}}}

% Entries of a matrix
\def\emLambda{{\Lambda}}
\def\emA{{A}}
\def\emB{{B}}
\def\emC{{C}}
\def\emD{{D}}
\def\emE{{E}}
\def\emF{{F}}
\def\emG{{G}}
\def\emH{{H}}
\def\emI{{I}}
\def\emJ{{J}}
\def\emK{{K}}
\def\emL{{L}}
\def\emM{{M}}
\def\emN{{N}}
\def\emO{{O}}
\def\emP{{P}}
\def\emQ{{Q}}
\def\emR{{R}}
\def\emS{{S}}
\def\emT{{T}}
\def\emU{{U}}
\def\emV{{V}}
\def\emW{{W}}
\def\emX{{X}}
\def\emY{{Y}}
\def\emZ{{Z}}
\def\emSigma{{\Sigma}}

% entries of a tensor
% Same font as tensor, without \bm wrapper
\newcommand{\etens}[1]{\mathsfit{#1}}
\def\etLambda{{\etens{\Lambda}}}
\def\etA{{\etens{A}}}
\def\etB{{\etens{B}}}
\def\etC{{\etens{C}}}
\def\etD{{\etens{D}}}
\def\etE{{\etens{E}}}
\def\etF{{\etens{F}}}
\def\etG{{\etens{G}}}
\def\etH{{\etens{H}}}
\def\etI{{\etens{I}}}
\def\etJ{{\etens{J}}}
\def\etK{{\etens{K}}}
\def\etL{{\etens{L}}}
\def\etM{{\etens{M}}}
\def\etN{{\etens{N}}}
\def\etO{{\etens{O}}}
\def\etP{{\etens{P}}}
\def\etQ{{\etens{Q}}}
\def\etR{{\etens{R}}}
\def\etS{{\etens{S}}}
\def\etT{{\etens{T}}}
\def\etU{{\etens{U}}}
\def\etV{{\etens{V}}}
\def\etW{{\etens{W}}}
\def\etX{{\etens{X}}}
\def\etY{{\etens{Y}}}
\def\etZ{{\etens{Z}}}

% The true underlying data generating distribution
\newcommand{\pdata}{p_{\rm{data}}}
% The empirical distribution defined by the training set
\newcommand{\ptrain}{\hat{p}_{\rm{data}}}
\newcommand{\Ptrain}{\hat{P}_{\rm{data}}}
% The model distribution
\newcommand{\pmodel}{p_{\rm{model}}}
\newcommand{\Pmodel}{P_{\rm{model}}}
\newcommand{\ptildemodel}{\tilde{p}_{\rm{model}}}
% Stochastic autoencoder distributions
\newcommand{\pencode}{p_{\rm{encoder}}}
\newcommand{\pdecode}{p_{\rm{decoder}}}
\newcommand{\precons}{p_{\rm{reconstruct}}}

\newcommand{\laplace}{\mathrm{Laplace}} % Laplace distribution

\newcommand{\E}{\mathbb{E}}
\newcommand{\Ls}{\mathcal{L}}
\newcommand{\R}{\mathbb{R}}
\newcommand{\emp}{\tilde{p}}
\newcommand{\lr}{\alpha}
\newcommand{\reg}{\lambda}
\newcommand{\rect}{\mathrm{rectifier}}
\newcommand{\softmax}{\mathrm{softmax}}
\newcommand{\sigmoid}{\sigma}
\newcommand{\softplus}{\zeta}
\newcommand{\KL}{D_{\mathrm{KL}}}
\newcommand{\Var}{\mathrm{Var}}
\newcommand{\standarderror}{\mathrm{SE}}
\newcommand{\Cov}{\mathrm{Cov}}
% Wolfram Mathworld says $L^2$ is for function spaces and $\ell^2$ is for vectors
% But then they seem to use $L^2$ for vectors throughout the site, and so does
% wikipedia.
\newcommand{\normlzero}{L^0}
\newcommand{\normlone}{L^1}
\newcommand{\normltwo}{L^2}
\newcommand{\normlp}{L^p}
\newcommand{\normmax}{L^\infty}

\newcommand{\parents}{Pa} % See usage in notation.tex. Chosen to match Daphne's book.

\DeclareMathOperator*{\argmax}{arg\,max}
\DeclareMathOperator*{\argmin}{arg\,min}

\DeclareMathOperator{\sign}{sign}
\DeclareMathOperator{\Tr}{Tr}
\let\ab\allowbreak


\usepackage{hyperref}
\usepackage{url}
\usepackage{booktabs}
\usepackage{graphicx}
\usepackage{caption}
\usepackage{booktabs}
\usepackage{multirow}
\usepackage{xcolor}
\usepackage{enumitem}
\usepackage[export]{adjustbox} % 必须带上 [export]

\definecolor{bestcyan}{HTML}{C1E1C1}
\definecolor{secondcyan}{RGB}{211,246,248}
\newcommand{\best}[1]{\cellcolor{bestcyan}\textbf{#1}}
\newcommand{\second}[1]{\cellcolor{secondcyan}#1}

\newcommand{\metricstd}[2]{#1\,{\scriptsize(#2)}}
\usepackage{graphicx}
\usepackage{wrapfig}
\usepackage{tabularx}
\usepackage{amsmath}
\usepackage{amssymb}
\usepackage{bm}
\usepackage{natbib}
\usepackage{algorithm}
\usepackage{algorithmic}

\title{PhyCAD: Physics-Aware Text-to-CAD Generation via Hierarchical Reinforcement Learning}

\author{
\textbf{Yuxin Cui},\,
\textbf{Weijie Zheng},\,
\textbf{Zhengyang Li},\,
\textbf{Han Wang},\,
\textbf{Wenjie Feng},\,
\textbf{Qiyu Kang}\thanks{Corresponding author: \texttt{qiyukang@ustc.edu.cn}} \\[0.5em]
University of Science and Technology of China, Hefei, China \\[0.3em]
}

\newcommand{\fix}{\marginpar{FIX}}
\newcommand{\new}{\marginpar{NEW}}

\iclrfinalcopy % Uncomment for camera-ready version, but NOT for submission.
\begin{document}

\maketitle
\fancyhead{}
\fancyhead[L]{Preprint}

\begin{figure}[h]
    \centering
    \includegraphics[width=\textwidth]{fig/teaser_cad.pdf}
    \caption{
\textbf{Physics-aware CAD optimization.}
\textbf{(A)} A source CAD model is annotated with FEM conditions and responses.
\textbf{(B)} Its exact dimensions are redacted to form an inverse-design task with physical constraints.
\textbf{(C)} While the SFT design achieves lower mass but violates the physical constraints, PhyCAD produces a feasible design with lower mass than the reference.
}
    \label{fig:teaser}
\end{figure}

\begin{abstract}
Language-driven computer-aided design (CAD) generation enables the creation of editable parametric models from text, yet geometric correctness alone does not ensure physical feasibility.
We present PhyCAD, a framework for physics-aware Text-to-CAD generation that trains language models to produce executable CadQuery programs under geometric and physical constraints.
To support this setting, we construct a Text-to-CAD dataset that links parametric CAD models with geometric specifications, loading conditions, and physical constraints.
We further develop an automated CAD--CAE verifier that combines program execution, geometric checks, and finite element analysis to provide feedback for model training.
A key learning challenge is to prioritize joint satisfaction of geometric and physical constraints while retaining informative feedback on partial progress.
A flat aggregation of heterogeneous rewards can obscure this priority, allowing improvements in individual criteria to compensate for unmet design requirements.
We therefore introduce a Hierarchical Balanced Advantage for Group Relative Policy Optimization, grouping candidates into levels according to their verification status and combining between-level progress with objective-wise peer ranking within each level.
This formulation preserves cross-level priority in the training signal while providing fine-grained feedback for satisfying constraints and improving feasible designs.
Experiments across multiple language models show that PhyCAD consistently improves joint geometric and physical constraint satisfaction.
These results demonstrate the value of verifiable engineering feedback for advancing language-driven CAD generation beyond geometric reconstruction toward physics-aware design. The code is available at \url{https://github.com/PhysAGI/PhyCAD}.

\end{abstract}

\section{Introduction}
% Computer-aided design (CAD) is a fundamental component of modern engineering, providing editable and parametric representations for downstream simulation, optimization, and manufacturing. Learning-based CAD generation has progressively moved from modeling CAD construction sequences and design histories~\citep{wu2021deepcad,willis2021fusion} toward language-driven generation with large language models (LLMs)~\citep{zhang2026large}. Recent methods can translate natural-language descriptions into parametric CAD sequences~\citep{khan2024text2cad}, generate CAD programs from textual or multimodal inputs~\citep{li2025llm4cad,sun2025large}, or directly synthesize executable modeling code through iterative or tool-augmented workflows~\citep{badagabettu2024query2cad,xie2025text,mallis2025cad}. In particular, CadQuery provides an executable, Python-based parametric representation that aligns naturally with the code-generation capabilities of pretrained LLMs~\citep{xie2025text,guan2026cad}.

% Despite this progress, reliably translating textual specifications into precise engineering structures remains challenging. Existing Text-to-CAD systems are commonly optimized or evaluated using geometric similarity, or visual consistency~\citep{khan2024text2cad,alrashedy2025generating,xie2025text,guan2026cad}. However, natural-language descriptions do not always establish an exact correspondence with the dimensions, topology, and structural details required by an engineering design, and small geometric deviations may substantially affect the behavior of the resulting structure. More importantly, engineering design is ultimately concerned with whether a component satisfies prescribed requirements under specific materials, loads, boundary conditions, and dimensional constraints, rather than merely whether it resembles a target shape. In conventional workflows, bringing an initially generated model into compliance often requires repeated manual measurement, geometric modification, simulation, and re-evaluation, resulting in a costly and labor-intensive refinement process. These limitations have motivated increasing attention to functional and physics-based criteria for evaluating and refining generated CAD models, allowing engineering feedback to guide geometry toward designs that satisfy downstream requirements~\citep{berger2026physics,son2026self}.

% Computer-aided engineering (CAE), particularly the finite element method (FEA), provides a natural source of deterministic feedback for evaluating generated designs. Recent studies have demonstrated the use of LLMs in automated FEA workflows and the acquisition of verifiable engineering signals~\citep{hou2025autofea,qi2025feagpt}. Meanwhile, reinforcement learning (RL) provides a natural mechanism for exploiting such automatically verifiable feedback~\citep{guan2026cad}, enabling simulation outcomes to be further incorporated into iterative CAD generation and optimization~\citep{berger2026physics,son2026self,deng2026cosmo}. Together, these approaches allow generated designs to be evaluated from both geometric and engineering perspectives. However, the resulting verification objectives are not independent, but exhibit a clear prerequisite structure: basic program validity and geometric compliance determine whether a generated design can be reliably carried forward to downstream engineering analysis, while physics-based evaluation becomes meaningful only after these lower-level requirements have been satisfied. Simply combining all objectives into a single flat reward may obscure these dependencies and make it difficult to effectively optimize objectives at different levels simultaneously. This motivates a hierarchical training strategy that progressively incorporates verification objectives according to their underlying dependency structure.

% In this work, we formulate physics-constrained Text-to-CAD generation, where a model receives textual geometric specifications together with materials, loading and boundary conditions, and physics-based constraints, and generates executable CadQuery programs. We construct a dataset for this setting and develop a complete CAD--CAE verification pipeline that automatically executes generated programs, verifies geometric and dimensional requirements, performs finite-element simulation, and converts deterministic verification outcomes into targeted reward signals. Building on supervised fine-tuning, we further introduce a hierarchical reinforcement-learning framework based on group-relative policy optimization~\citep{shao2024deepseekmath}, in which verification objectives are progressively incorporated according to their dependency structure. This enables physics-based verification to directly participate in model optimization, providing a controlled setting for studying how geometric and physical feedback can jointly guide Text-to-CAD generation.

CAD is a fundamental component of modern engineering, providing editable and parametric representations for downstream simulation, optimization, and manufacturing. Learning-based CAD generation has progressively moved from modeling CAD construction sequences and design histories~\citep{wu2021deepcad,willis2021fusion} toward language-driven generation with large language models (LLMs)~\citep{zhang2026large}. Recent methods can translate natural-language descriptions into parametric CAD sequences~\citep{khan2024text2cad}, generate CAD programs from textual or multimodal inputs~\citep{li2025llm4cad,sun2025large}, or synthesize executable modeling code through iterative or tool-augmented workflows~\citep{badagabettu2024query2cad,xie2025text,mallis2025cad}. In particular, CadQuery provides an executable, Python-based parametric representation that aligns naturally with code-generation capabilities of pretrained LLMs~\citep{xie2025text,guan2026cad}.

Despite this progress, geometric correctness alone does not ensure physical feasibility. Existing Text-to-CAD systems are commonly optimized or evaluated using geometric similarity or visual consistency~\citep{khan2024text2cad,alrashedy2025generating,xie2025text,guan2026cad}. 
However, natural-language descriptions do not always fully specify the dimensions, topology, and structural details required for engineering design, and even small geometric deviations may substantially affect structural behavior. More importantly, engineering design is ultimately concerned with whether a component meets prescribed dimensional and physical requirements under specified material properties, loads, and boundary conditions, rather than merely whether it resembles a target shape.
Bringing an initially generated model into compliance often requires repeated manual measurement, geometric modification, simulation, and re-evaluation, resulting in a costly refinement process. These limitations have motivated increasing attention to functional and physics-based criteria for evaluating and refining generated CAD models, allowing engineering feedback to guide geometry toward designs that satisfy downstream requirements~\citep{berger2026physics,son2026self}.

Computer-aided engineering (CAE), particularly simulation based on the finite element method (FEM), provides a natural source of deterministic feedback for evaluating generated designs. Recent studies have demonstrated the use of LLMs in automated FEM-based workflows and the acquisition of verifiable engineering signals~\citep{hou2025autofea,qi2025feagpt}. Simulation outcomes have also been used to guide iterative CAD generation and optimization~\citep{berger2026physics,son2026self,deng2026cosmo}. Reinforcement learning (RL) has further enabled CAD-generation models to learn from geometric feedback~\citep{guan2026cad}. Training with both geometric and physical feedback introduces a key challenge: prioritizing joint constraint satisfaction while retaining informative feedback on partial progress. A flat aggregation of heterogeneous rewards can obscure this priority, allowing improvements in individual criteria to compensate for unmet design requirements. This calls for a learning signal that distinguishes progress toward feasibility from improvements within the feasible design space.

As illustrated in Fig.~\ref{fig:teaser}, we present PhyCAD, a physics-aware Text-to-CAD framework that generates executable CadQuery programs from geometric specifications and physical constraints under prescribed material properties, loads, and boundary conditions. The task defines a constrained design space with admissible dimensional ranges rather than a unique target geometry. We construct a corresponding dataset and an automated CAD--CAE verifier that provides training feedback through program execution, geometric checks, and FEM-based simulation. Building on supervised fine-tuning, we introduce a Hierarchical Balanced Advantage for Group Relative Policy Optimization (GRPO)~\citep{shao2024deepseekmath}, grouping candidates by verification status and combining between-level progress with objective-wise peer ranking within each level. This formulation preserves cross-level priority in the training signal while providing fine-grained feedback for satisfying constraints and improving feasible designs. The trained model directly generates programs without an iterative simulation-guided refinement loop at inference time. Experiments across multiple language models show that PhyCAD consistently improves joint geometric and physical constraint satisfaction.


Our contributions are summarized as follows:
\begin{itemize}
[leftmargin=*, itemsep=0.15em, topsep=0.1em, parsep=0pt, partopsep=0pt]
\item We formulate physics-aware Text-to-CAD generation and construct a dataset that jointly captures geometric specifications, loading and boundary conditions, and physics-based constraints.
\item We develop an automated CAD--CAE verification pipeline that provides verifiable feedback on program executability, geometric and dimensional compliance, and physical constraint satisfaction to support model training.
\item We introduce a Hierarchical Balanced Advantage that combines between-level progress with objective-wise peer ranking within each level, preserving cross-level priority in the training signal while providing fine-grained feedback for satisfying constraints and improving feasible designs.
\end{itemize}



% Despite this progress, geometric correctness alone is insufficient for engineering design. Existing Text-to-CAD systems are commonly optimized or evaluated using program validity, geometric similarity, or visual consistency~\citep{khan2024text2cad,alrashedy2025generating,xie2025text,guan2026cad}. These criteria measure whether a generated program executes or reproduces the intended shape, but do not directly determine whether the resulting artifact satisfies engineering requirements under prescribed materials, loads, boundary conditions, and dimensional constraints. A geometrically plausible design may still exhibit excessive displacement, high stress, insufficient stiffness, or invalid load transfer. Recent studies therefore argue for evaluating generative CAD using functional and physics-based criteria rather than geometry alone~\citep{berger2026physics,son2026self}.

% This has motivated the integration of computer-aided engineering (CAE), particularly the finite element method (FEA), into LLM-based engineering workflows. AutoFEA and FeaGPT demonstrate that LLMs can assist or orchestrate FEA configuration, geometry processing, meshing, simulation, and result analysis~\citep{hou2025autofea,qi2025feagpt}. Physics-in-the-loop CAD agents further use deterministic simulation feedback to iteratively revise generated designs~\citep{berger2026physics,son2026self}, while COSMO-Agent formulates closed-loop CAD--CAE optimization as a tool-augmented reinforcement-learning problem with multi-constraint rewards~\citep{deng2026cosmo}. Meanwhile, CAD-Coder shows that reinforcement learning can improve Text-to-CAD generation using geometric and format rewards~\citep{guan2026cad}. These advances establish the value of verifiable feedback, but leave open how to organize heterogeneous CAD--CAE signals when they exhibit clear prerequisite dependencies: a candidate must first be executable, then geometrically compliant, before its physical performance can be meaningfully evaluated.



\section{Related Work}
% \subsection{Text-to-CAD Generation}
% \paragraph{Text-to-CAD Generation.}
% Learning-based CAD generation has evolved from structured design modeling to language-conditioned generation with pretrained models. Early work modeled CAD construction histories, sketch constraints, and assembly relations using datasets such as DeepCAD~\citep{wu2021deepcad}, SketchGraphs~\citep{seff2020sketchgraphs}, and Fusion 360 Gallery~\citep{willis2021fusion,willis2022joinable}. Task-specific transformers have subsequently been developed to map text or engineering drawings to CAD command sequences~\citep{khan2024text2cad,qin2025drawing2cad}. Pretrained LLMs and VLMs have also been adapted to generate code-like or structured CAD representations~\citep{li2025cad,govindarajan2025cadmium,yuan2024openecad}. Programmatic representations such as CadQuery support directly executable generation~\citep{xie2025text}, while tool-augmented, self-refinement, and evolutionary approaches improve generated CAD programs through iterative feedback~\citep{mallis2025cad,li2026seek,preintner2025evocad}. Our work studies Text-to-CAD generation under joint geometric and physical requirements, including global dimensional bounds and FEM-verified physical constraints.

% \paragraph{Physics-Aware CAD Generation and Engineering Agents.}
% Recent work has started to evaluate generated CAD models against engineering requirements rather than geometric similarity alone. \citep{berger2026physics} integrate physics-based verification into a multi-agent CAD design and refinement loop. \citep{son2026self} study CAD generation from single-part and multi-part engineering briefs, using structured blueprints, visual inspection, and finite-element feedback to iteratively refine generated STEP models. COSMO-Agent~\citep{deng2026cosmo} uses tool-augmented reinforcement learning to train an LLM for closed-loop CAD--CAE optimization, where the agent updates geometric parameters and material choices for parameterized components according to simulation feedback. Related LLM-based systems automate workflows in computational fluid dynamics~\citep{chen2024metaopenfoam,xu2025cfdagent,yue2025foam}, finite-element analysis~\citep{hou2025autofea,qi2025feagpt}, and multiphysics simulation~\citep{zhang2025mooseagent}. 
% Complementary to iterative CAD--CAE refinement, PhyCAD focuses on improving the engineering feasibility of designs at generation time. By learning from CAD--CAE verification outcomes across design tasks, the model is explicitly optimized to jointly satisfy geometric specifications and physical constraints when generating executable CadQuery programs. Engineering feedback thus improves the generator's ability to produce compliant designs, without requiring a simulation-guided refinement loop at inference time.

% \paragraph{Reinforcement Learning with Verifiable Feedback.}
% RL with verifiable feedback provides a practical way to optimize language models when correctness can be evaluated automatically. GRPO~\citep{shao2024deepseekmath} enables policy optimization using group-relative rewards, while InternBootcamp~\citep{li2025internbootcamp} provides diverse tasks with programmatic verifiers. In CAD, CAD-Coder~\citep{guan2026cad} combines supervised fine-tuning with GRPO using geometric and format rewards for Text-to-CAD generation, while cadrille~\citep{kolodiazhnyi2025cadrille} uses geometric feedback for multimodal CAD reconstruction. ReCAD~\citep{li2026recad} combines parameterized-code guidance with hierarchical primitive learning, organizing training over increasingly complex geometric structures. These methods primarily optimize geometric fidelity and output validity rather than FEM-verified physical constraint satisfaction. PhyCAD instead incorporates CAD--CAE feedback into Text-to-CAD policy training through a Hierarchical Balanced Advantage. Unlike a curriculum over geometric primitives, our hierarchy organizes candidates by verification status, combining between-level progress with objective-wise peer ranking within each level.


\paragraph{Text-to-CAD Generation.}
Early learning-based CAD methods modeled construction histories, sketch constraints, and assembly relations using datasets such as DeepCAD~\citep{wu2021deepcad}, SketchGraphs~\citep{seff2020sketchgraphs}, and Fusion 360 Gallery~\citep{willis2021fusion,willis2022joinable}. Subsequent approaches map text or engineering drawings to CAD sequences using task-specific transformers~\citep{khan2024text2cad,qin2025drawing2cad}, or adapt pretrained LLMs and VLMs to structured CAD generation~\citep{li2025cad,govindarajan2025cadmium,yuan2024openecad}. CadQuery supports executable program generation~\citep{xie2025text}, while tool augmentation, self-refinement, and evolutionary search enable iterative refinement~\citep{mallis2025cad,li2026seek,preintner2025evocad}. PhyCAD studies generation under joint geometric and physical requirements, including global dimensional bounds and FEM-verified physical constraints.

\paragraph{Physics-Aware CAD Generation and Engineering Agents.}
Recently, \citep{berger2026physics} integrates physics-based verification into multi-agent CAD refinement, while \citep{son2026self} combines structured blueprints, visual inspection, and finite-element feedback to refine single-part and multi-part STEP models. COSMO-Agent~\citep{deng2026cosmo} trains an LLM through tool-augmented RL to optimize parameterized components through CAE-guided updates to geometry and material choices. Related agents automate computational fluid dynamics~\citep{chen2024metaopenfoam,xu2025cfdagent,yue2025foam}, finite-element analysis~\citep{hou2025autofea,qi2025feagpt}, and multiphysics simulation~\citep{zhang2025mooseagent}. Complementary to these engineering workflows, PhyCAD addresses a learning challenge in physics-aware CAD generation: prioritizing joint geometric and physical feasibility while retaining informative feedback on partial progress. We introduce a Hierarchical Balanced Advantage that combines progress across verification levels with objective-wise peer ranking within each level. This organizes heterogeneous CAD--CAE outcomes into a training signal that preserves cross-level priority while guiding constraint satisfaction and improvements within the feasible design space.

\paragraph{Reinforcement Learning with Verifiable Feedback.}
GRPO~\citep{shao2024deepseekmath} enables policy optimization using group-relative rewards, while InternBootcamp~\citep{li2025internbootcamp} provides diverse tasks with programmatic verifiers. In CAD, CAD-Coder~\citep{guan2026cad} combines supervised fine-tuning and GRPO using geometric and format rewards, while cadrille~\citep{kolodiazhnyi2025cadrille} uses geometric feedback for multimodal reconstruction. ReCAD~\citep{li2026recad} combines parameterized-code guidance with hierarchical primitive learning, following a curriculum over increasingly complex geometric structures. These methods primarily target geometric fidelity and output validity rather than FEM-verified physical constraints. PhyCAD incorporates CAD--CAE feedback into policy training through a Hierarchical Balanced Advantage. Rather than scheduling a geometric curriculum, it groups candidates by verification status and combines between-level progress with objective-wise peer ranking within each level.

\section{Method}
\begin{figure}[h]
    \centering
    \includegraphics[width=\textwidth]{fig/framework.pdf}
    \caption{
    \textbf{Overview of our framework.}
    The LLM policy generates candidate CAD programs from a physics-conditioned prompt containing geometric requirements, design-domain constraints, loads, and physical limits.
    A deterministic CAD-CAE verifier sequentially evaluates each program for {CAD executability}, {geometric validity}, and {FEM-based physical feasibility and objectives}.
    These verification results define a {Hierarchical Balanced Advantage}, which combines between-level progress with within-level peer ranking over objective scores.
    The resulting advantage $A_i$ provides structured reinforcement signals for updating the CAD-generation policy.
    }
    \label{fig:overview}
    \vspace{-1em}
\end{figure}

Our goal is to train a LLM to generate CAD programs that satisfy both geometric specifications and physical constraints. As illustrated in Fig.~\ref{fig:overview}, our framework takes a physics-aware prompt as input, generates candidate CadQuery programs with an LLM policy, evaluates them through a deterministic CAD--CAE verification pipeline, and converts the verification outcomes into structured reinforcement signals for policy optimization. The resulting training framework integrates CAD executability, geometric compliance, and physics-based feasibility into a unified optimization process.

% \subsection{Problem Formulation} 
% \label{sec:problem_formulation}

% Let \(q\in\mathcal{Q}\) denote a physics-aware design prompt, where \(\mathcal{Q}\) is the space of design specifications. Each prompt jointly describes geometric requirements, the admissible design domain, material properties, loading and boundary conditions, engineering constraints, and a design objective. An LLM policy \(\pi_{\theta}\), parameterized by \(\theta\in\mathbb{R}^{d}\), defines a conditional distribution over the space of CadQuery programs \(\mathcal{O}\). We sample a group of candidate programs as \(o_i\sim\pi_{\theta}(\cdot\mid q)\), where \(o_i\in\mathcal{O}\) and \(i=1,\ldots,G\).

% Our setting specifies selected geometric parameters through admissible ranges rather than fixed values, defining a constrained design space rather than a unique target geometry. A single prompt may therefore admit multiple geometrically distinct solutions. 
% A program \(o\in\mathcal O\) is feasible if it executes successfully to produce a valid CAD solid, its geometry satisfies the specified requirements, and its FEM evaluation succeeds with all physical constraints satisfied. We denote the set of feasible programs for prompt \(q\) by \(\mathcal F(q)\subseteq\mathcal O\).
% Feasibility alone does not select a unique design: different dimensional choices may satisfy the same requirements while producing different structural masses. We therefore use an engineering objective \(J(o;q)\) to express the preference among feasible programs, with smaller values preferred. For the tasks considered here, \(J(o;q)=m(o;q)\), where \(m(o;q)\) is the mass of the generated solid under the specified material properties. The resulting constrained design problem is
% \begin{align}
%      o^{*}\in\underset{o\in\mathcal F(q)}{\arg\min}\;J(o;q).
%     \label{eq:problem}
% \end{align}
% The prescribed mass upper bound contributes to feasibility, whereas the objective favors further mass reduction among feasible designs. The training feedback introduced below additionally considers capped stress and displacement margins when comparing feasible candidates.


% During training, optimization is performed over the sampled candidate set
% \(\mathcal{O}_{\theta}(q)\). The CAD--CAE verifier determines which candidates
% belong to the feasible subset \(\mathcal{F}_{\theta}(q)\) and provides
% hierarchical feedback for infeasible candidates according to their verification
% status. Accordingly, we optimize \(\pi_{\theta}\) to increase the likelihood
% of feasible designs while improving the engineering objective \(J(o;q)\)
% among feasible candidates.
\subsection{Problem Formulation} 
\label{sec:problem_formulation}

Let \(q\in\mathcal{Q}\) denote a physics-aware design prompt describing geometric requirements, the admissible design domain, material properties, loading and boundary conditions, engineering constraints, and a design objective. An LLM policy \(\pi_{\theta}\), parameterized by \(\theta\in\mathbb{R}^{d}\), defines a conditional distribution over CadQuery programs \(\mathcal{O}\). During training, we sample a group of candidates
\(o_i\sim\pi_{\theta}(\cdot\mid q)\), \(i=1,\ldots,G\).

Our setting specifies selected geometric parameters through admissible ranges rather than fixed values, defining a constrained design space rather than a unique target geometry. A program \(o\in\mathcal{O}\) is feasible if it executes successfully to produce a valid CAD solid, satisfies the specified geometric requirements, and passes FEM evaluation with all physical constraints satisfied. We denote the feasible set by \(\mathcal{F}(q)\subseteq\mathcal{O}\). Among feasible programs, we optimize an engineering objective \(J(o;q)\), with smaller values preferred. For the tasks considered here, \(J(o;q)=m(o;q)\), where \(m(o;q)\) is the mass of the generated solid. The resulting constrained design problem is
\begin{align}
    o^{*}\in\underset{o\in\mathcal F(q)}{\arg\min}\;J(o;q).
    \label{eq:problem}
\end{align}
During training, optimization is performed over the sampled candidate set
\(\mathcal{O}_{\theta}(q)\). The CAD--CAE verifier identifies the feasible subset
\(\mathcal{F}_{\theta}(q)\) and provides hierarchical feedback for candidates
according to their verification status. The policy is optimized to increase the
probability of feasible designs while improving \(J(o;q)\) among feasible candidates.

% Our goal is to train a language model to generate CAD programs that satisfy both geometric specifications and engineering requirements. As illustrated in Fig.~\ref{fig:overview}, our framework takes a physics-conditioned prompt as input, generates candidate CadQuery programs with an LLM policy, evaluates them through a deterministic CAD--CAE verification pipeline, and converts the verification outcomes into structured reinforcement signals for policy optimization. The resulting training framework integrates CAD executability, geometric compliance, and physics-based feasibility into a unified optimization process.

% \subsection{Problem Formulation}
% \label{sec:problem_formulation}

% Let \(q\in\mathcal{Q}\) denote a physics-aware design prompt, where
% \(\mathcal{Q}\) is the space of design specifications.
% Each prompt jointly describes the target geometry, admissible design
% domain, material properties, loading and boundary conditions,
% engineering constraints, and design objective.
% An LLM policy
% \(\pi_{\theta}\), parameterized by
% \(\theta\in\mathbb{R}^{d}\), defines a conditional distribution over
% the space of CadQuery programs \(\mathcal{O}\).
% We sample a group of candidate programs as
% \(o_i \sim \pi_{\theta}(\cdot\mid q)\),
% where \(o_i\in\mathcal{O}\) and \(i=1,\ldots,G\).

% Unlike conventional Text-to-CAD generation, which typically aims to reproduce a largely predetermined geometry from explicit dimensional specifications, our setting relaxes selected geometric parameters from fixed values to admissible ranges, thereby defining a constrained design space. The model is required to both instantiate a valid CAD geometry within this space and identify a configuration that satisfies the prescribed physics-based constraints under the given engineering conditions. Specifically, a candidate program \(o_i\) is considered
% feasible only if it can be successfully executed into a valid CAD
% solid, the resulting geometry satisfies the specified geometric
% requirements, and its physical response under FEM
% satisfies the prescribed engineering limits. 
% These requirements are
% ordered by construction: geometric verification is meaningful only
% for an executable CAD model, while physical verification is meaningful
% only for a geometrically valid design.

% Given a prompt \(q\), the policy \(\pi_{\theta}\) generates a group of
% \(G\) candidate programs,
% \(
% \mathcal{O}_{\theta}(q)=\{o_i\}_{i=1}^{G}
% \).
% Let \(\mathcal{F}(q)\subseteq\mathcal{O}\) denote the set of all
% CadQuery programs that satisfy the CAD, geometric, and physical
% requirements specified by \(q\).
% The feasible candidates produced by the current policy are therefore
% given by
% \(
% \mathcal{F}_{\theta}(q)
% =
% \mathcal{O}_{\theta}(q)\cap\mathcal{F}(q)
% \).
% For each feasible program \(o\in\mathcal{F}(q)\), let
% \(J(o;q)\in\mathbb{R}\) denote the corresponding engineering objective,
% such as structural mass. The underlying engineering-constrained Text-to-CAD problem is formulated as
% \begin{equation}
%     o^{*}
%     \in
%     \underset{o\in\mathcal{F}(q)}{\arg\min}
%     \; J(o;q).
%     \label{eq:problem}
% \end{equation}
% During training, optimization is performed over the sampled candidate set
% \(\mathcal{O}_{\theta}(q)\). The CAD--CAE verifier determines which candidates
% belong to the feasible subset \(\mathcal{F}_{\theta}(q)\) and provides
% hierarchical feedback for infeasible candidates according to their verification
% status. Accordingly, we optimize \(\pi_{\theta}\) to shift probability mass
% toward feasible designs while improving the engineering objective \(J(o;q)\)
% within \(\mathcal{F}_{\theta}(q)\). 



\subsection{Deterministic CAD--CAE Verifier}
\label{sec:verifier}

For any group of programs \(o_i\) sampled from \(\pi_{\theta}(\cdot \mid q)\), we evaluate each program using the same deterministic CAD--CAE pipeline. The verifier proceeds through three stages.

\paragraph{CAD execution.}
The generated program is checked for syntax and safety, executed, and tested for valid output construction. It must produce a valid CAD solid that can be exported in STEP format and reimported by the downstream geometry kernel. Failure at this stage is recorded as a CAD-construction failure.

\paragraph{Geometry check.}
For successfully constructed solids, we compare the global bounding-box extents with the prompt's dimensional ranges and check whether the required geometric constructions are present. This checks compliance with the design requirements, separately from the solid-validity check above. Candidates satisfying these requirements proceed to CAE verification.

\paragraph{CAE verification.}
The solid is meshed and analyzed using FEM under the prescribed material properties, loads, and boundary conditions. For candidate \(o_i\), we record the volume-weighted 99th-percentile element von Mises stress \(\sigma_i\), maximum nodal displacement \(\delta_i\), and mass \(m_i=m(o_i;q)\). Let \(\bar\sigma\), \(\bar\delta\), and \(\bar m\) denote their respective upper bounds in the prompt. Physical verification succeeds when the CAE procedure completes and \(\sigma_i\leq\bar\sigma\), \(\delta_i\leq\bar\delta\), and \(m_i\leq\bar m\).

To summarize these outcomes, let \(L_i\in\{0,1,2,3\}\) denote the verification level of candidate \(o_i\):

\begin{equation}
L_i=
\begin{cases}
0, & \text{CAD construction fails},\\
1, & \text{CAD succeeds but geometric requirements fail},\\
2, & \text{geometry passes but CAE or a physical constraint fails},\\
3, & \text{all geometric and physical requirements pass}.
\end{cases}
\label{eq:level}
\end{equation}

Thus, the three checks yield four possible outcomes: failure at any one of the checks, or success at all three. A larger \(L_i\) indicates further progress through the verification procedure; \(L_i=3\) identifies a jointly feasible candidate. %The individual measurements provide finer distinctions within each level.

\subsection{Hierarchical Balanced Advantage} \label{sec:hba} 
The verification levels distinguish coarse progress toward feasibility but do not distinguish candidates with the same status. For example, two physically infeasible candidates may violate their constraints by very different amounts. Conversely, directly aggregating all objective scores can allow improvements in individual criteria to compensate for unmet requirements. We therefore combine a between-level advantage with a within-level comparison, retaining both the verification ordering and feedback on partial progress.

\paragraph{Between levels.} For \(G\) rollouts generated from the same prompt, we define the level advantage as 
\begin{equation} 
A_i^{\mathrm{level}} = L_i - \frac{1}{G} \sum_{j=1}^{G} L_j. 
\label{eq:level_advantage} 
\end{equation} 
This term compares each candidate's level with the group mean. Candidates above the mean receive positive level advantages. When all candidates occupy the same level, this term is zero, and the within-level component provides the remaining distinctions.

\begin{figure}[t]
    \centering
    \includegraphics[width=\textwidth]{fig/reward_new.pdf}
    % \caption{
    % \textbf{Rewards for within-level PeerRank.}
    % Each level evaluates a different aspect, from CAD executability (\(L_0\)) and geometric compliance (\(L_1\)) to physical feasibility (\(L_2\)) and feasible-design quality (\(L_3\)). The blocks show the objective definitions and their relative weights within each level. Each objective is independently peer-ranked among rollouts in the same level, and the resulting centered ranks are combined using these weights to form the within-level advantage \(\widetilde{A}_i\).
    % }
    \caption{
\textbf{Reward definitions for within-level PeerRank.}
Each block specifies how an evaluated outcome is mapped to an objective reward
within its corresponding level.
Solid regions denote fixed rewards once a condition is satisfied, while shaded
gradients indicate rewards that vary continuously with closeness to the target
or the achieved margin; faded regions correspond to outcomes receiving zero or
lower reward.
The resulting objective rewards are independently peer-ranked and combined
using the indicated within-level weights.
}
    \label{fig:reward}
    \vspace{-1em}
\end{figure}

\paragraph{Within a level.}
Candidates are compared only with others from the same prompt and the same verification level. We first map the verification outcomes to level-specific objective scores, then rank candidates separately for each objective. Fig.~\ref{fig:reward} summarizes the score definitions and weights.
\textbf{(1).} For candidates with \(L_i=0\), CAD-progress feedback reflects output completeness, code-only formatting, syntax validity, safe calls, result construction, and STEP export. This distinguishes partial construction progress even when a valid CAD solid has not been obtained.
\textbf{(2).}  For \(L_i=1\), we assign equal weights to bounding-box closeness and geometric-semantic recall. These scores measure proximity to the dimensional ranges and coverage of the required geometric constructions, respectively.
\textbf{(3).}  For \(L_i=2\), we compare closeness to the stress, displacement, and mass limits, again with equal weights. For an available measurement \(v\) with upper bound \(\bar v\), its score is \(1\) when \(v\leq\bar v\), and \(\bar v/v\) otherwise. An unavailable measurement receives a score of zero. These scores distinguish different degrees of constraint violation without rewarding further improvement of a quantity that already satisfies its limit.
\textbf{(4).}  For \(L_i=3\), all constraints are satisfied. We compare mass efficiency \(1-m_i/\bar m\), stress margin, and displacement margin with equal weights. The two margin scores are capped at a \(10\%\) margin, beyond which additional margin receives no further reward. Feasible-design feedback therefore combines the mass objective with capped margin objectives rather than ranking feasible candidates by mass alone.


For a given prompt and a level \(\ell\in\{0,1,2,3\}\), let \(\mathcal I_\ell=\{j\in\{1,\ldots,G\}:L_j=\ell\}\) be the set of candidate indices at that level. Let \(K_\ell\) be its number of objectives. For candidate \(i\in\mathcal I_\ell\), denote its score for objective \(k\in\{1,\ldots,K_\ell\}\) by \(s_{ik}\), and its tie-aware centered rank among \(\mathcal I_\ell\) by \(r_{ik}\in[-1,1]\), with higher ranks indicating better scores. Scores within \(10^{-3}\) are treated as ties. With nonnegative objective weights \(w_{\ell k}\) satisfying \(\sum_{k=1}^{K_\ell}w_{\ell k}=1\), the within-level advantage is

\begin{equation}
\widetilde A_i=
\begin{cases}
\displaystyle\sum_{k=1}^{K_{L_i}}w_{L_i k}\,r_{ik},
& |\mathcal I_{L_i}|>1,\\[2mm]
0, & |\mathcal I_{L_i}|=1.
\end{cases}
\label{eq:peer_rank}
\end{equation}

We refer to this objective-wise ranking and averaging operation as PeerRank. It compares ranks rather than a weighted sum of the original objective scores, and gives \(|\widetilde A_i|\leq1\).

% REVIEW RANK: The source specifies the range and tie tolerance but not the
% exact centered-rank transform or tolerance-based tie-grouping algorithm.
% The equation above makes the stated weighted average explicit; it does not
% complete the missing ranking definition. Insert the implementation-matched
% definition here or in an appendix. Do not assume midranks or pairwise ranks.

\paragraph{Combined advantage.}

The nonnegative weights \(\lambda_L\) and \(\lambda_w\) control the between-level and within-level terms, respectively. We combine them as
\begin{equation}
A_i=\lambda_L A_i^{\mathrm{level}}+\lambda_w\widetilde A_i.
\label{eq:hba}
\end{equation}

For candidates \(i\) and \(j\) from the same prompt group with \(L_i>L_j\), the common level baseline cancels, giving
\[
A_i-A_j
=\lambda_L(L_i-L_j)+\lambda_w(\widetilde A_i-\widetilde A_j)
\geq\lambda_L-2\lambda_w.
\]
Choosing \(\lambda_L>2\lambda_w\) therefore ensures that the higher-level candidate receives a larger advantage regardless of its within-level rank. This is an ordering guarantee for the training signal within each prompt group. Both components are computed together, not introduced in successive training stages.

\paragraph{Policy optimization.}
\label{sec:policy_optimization}

We use the hierarchical advantages in GRPO~\citep{shao2024deepseekmath}. The policy loss uses DAPO-style token-level aggregation: generated-token losses are averaged across the batch rather than first averaged separately within each completion. Let \(\mathbf A\) collect the advantages for the sampled completions, each computed within its own prompt group, and let \(\mathcal L_{\mathrm{DAPO}}(\theta;\mathbf A)\) denote the GRPO policy term with this aggregation. With a reference policy \(\pi_{\mathrm{ref}}\) and KL regularization weight \(\beta\geq0\), the training objective is

\begin{equation}
\mathcal L
=\mathcal L_{\mathrm{DAPO}}(\theta;\mathbf A)
+\beta D_{\mathrm{KL}}(\pi_\theta\,\|\,\pi_{\mathrm{ref}}),
\label{eq:grpo}
\end{equation}

where \(D_{\mathrm{KL}}\) denotes the Kullback--Leibler divergence. The CAD kernel and FEM solver supply verification outcomes, not gradients through geometry construction or simulation. At inference time, the trained policy directly generates a CadQuery program without an iterative simulation-guided refinement loop.



















% \paragraph{Within a level.}
% Candidates sharing the same verification level are compared only with their same-level peers. Each level uses a different set of objectives, all normalized to \([0,1]\).

% For \(L_0\), we use a CAD progress score that combines EOS termination, code-only formatting, syntax validity, safety, result-variable construction, and successful STEP export. This provides a graded signal even when a valid CAD solid has not yet been produced.

% For \(L_1\), candidates are compared using bounding-box closeness and geometric-semantic recall with equal weights. The former measures how closely the realized dimensions satisfy the prescribed range, while the latter measures whether all geometric capabilities required by the prompt are present.

% For \(L_2\), we compare physical constraint closeness for stress, displacement, and mass, again with equal weights. For each quantity \(v\) with limit \(\bar{v}\), the score is \(1\) if \(v \leq \bar{v}\), \(\bar{v}/v\) otherwise, and \(0\) when the corresponding CAE result is unavailable. Therefore, a failed FEM evaluation receives zero stress and displacement scores, while the mass score can still be recovered from the CAD geometry.

% For \(L3\), the design is already feasible, and the objectives shift from satisfying constraints to improving the feasible solution. We use mass efficiency \(1-m/\bar{m}\), together with stress and displacement margins. The latter two are capped at a \(10\%\) safety margin, so additional margin beyond this point provides no further benefit. As a result, mass reduction becomes the primary optimization objective once feasibility has been achieved.

% Within each level, every objective is independently converted into a tie-aware centered rank in \([-1,1]\), with values within \(10^{-3}\) treated as ties. The ranked objectives are then averaged using the corresponding level weights to obtain
% \begin{equation}
%     \widetilde{A}_i
%     =
%     \operatorname{PeerRank}(o_i \mid L_i).
%     \label{eq:peer_rank}
% \end{equation}
% If no same-level peer exists, we set \(\widetilde{A}_i=0\). Since ranking is performed separately for each objective, heterogeneous quantities such as geometry, stress, displacement, and mass do not need to be placed on a common numerical scale.

% \paragraph{Combined advantage.}
% The final advantage used for policy optimization is
% \begin{equation}
%     A_i
%     =
%     \lambda_L
%     \left(
%         L_i-\frac{1}{G}\sum_{j=1}^{G}L_j
%     \right)
%     +
%     \lambda_w \widetilde{A}_i,
%     \label{eq:hba}
% \end{equation}
% where \(\lambda_L=1\) and \(\lambda_w=0.25\). Because \(|\widetilde{A}_i|\leq 1\) and \(\lambda_L>2\lambda_w\), a candidate at a higher verification level always receives a larger advantage than a candidate at a lower level. The within-level term therefore refines the ordering without overriding the CAD--geometry--physics hierarchy.

% For monitoring, we additionally report a scalar reward defined as the verification level plus the weighted objectives of that level, resulting in four non-overlapping reward intervals \(L0:[0,1]\), \(L1:[1,2]\), \(L2:[2,3]\), and \(L3:[3,4]\). This scalar is used only for logging and does not directly determine the policy gradient.


% \subsection{Policy Optimization}
% \label{sec:policy_optimization}

% The hierarchical advantage \(A_i\) is used to optimize the CAD-generation policy with the DAPO objective,
% \begin{equation}
%     \mathcal{L}
%     =
%     \mathcal{L}_{\mathrm{DAPO}}(\theta; A_i)
%     +
%     \beta
%     D_{\mathrm{KL}}
%     \left(
%         \pi_{\theta}
%         \,\Vert\,
%         \pi_{\mathrm{ref}}
%     \right),
%     \label{eq:policy_update}
% \end{equation}
% where \(\pi_{\mathrm{ref}}\) denotes the reference policy and \(\beta\) controls the KL regularization strength.

% The CAD kernel and FEM solver are used only as deterministic external verifiers and therefore need not be differentiable. They convert each complete program rollout into a verification level and level-specific objectives, while the hierarchical advantage determines the actual policy update. This allows the model to improve according to the ordered structure of the task: first producing valid CAD, then satisfying geometric requirements, then achieving physical feasibility, and finally improving engineering performance within the feasible region.














\section{Experiment}
\begin{figure}[t]
    \centering
    \includegraphics[width=\textwidth]{fig/result_new.pdf}
\caption{
\textbf{Qualitative comparison across training stages and model families.}
Surface color indicates stress relative to the prescribed limit, while the bars
summarize displacement, stress, and mass constraint satisfaction.
Across model families, SFT improves CAD generation but can still violate geometric
or physical constraints, whereas GRPO more consistently achieves joint feasibility.
}
    \label{fig:qualitative_comparison}
    \vspace{-1em}
\end{figure}


\subsection{Dataset Generation}
\label{sec:dataset}

We construct a physics-aware Text-to-CAD dataset from the
Text-to-CadQuery corpus~\citep{xie2025text}, which contains approximately
\(110\text{k}\) CadQuery programs with geometric and semantic descriptions.
Each program serves as a reference design and is processed by a CAD--CAE
pipeline. We first derive boundary and loading conditions from the geometry,
then perform FEM to obtain the maximum displacement \(\delta_{\max}\),
\(99\)-th percentile von Mises stress \(\sigma_{\mathrm{P99}}\), and mass \(m\).

Each successfully analyzed reference is converted into an inverse-design task
rather than used as a direct reconstruction target. We remove exact dimensions
and other numerical specifications that could reveal the reference geometry,
while retaining high-level geometric semantics needed to define the intended
structure. The reference measurements are then used to derive displacement,
stress, and mass limits with a small feasibility margin, ensuring that the
reference remains a feasible solution without exposing its exact simulation
results. The resulting prompt specifies the geometry semantics, coarse design
bounds, boundary and loading conditions, and physical requirements, and asks
the model to generate a feasible design while minimizing mass. Detailed
construction and sanitization rules are provided in the Appendix.

After CAD--CAE annotation, \(77{,}067\) designs successfully pass the complete
pipeline. We divide them into disjoint SFT and GRPO subsets. The SFT subset contains \(11{,}024\) samples paired with sanitized
reference CadQuery programs, whereas the GRPO subset contains \(66{,}043\)
tasks without target programs and is optimized solely through CAD--CAE
verification.
% We construct a physics-aware Text-to-CAD dataset from the existing Text-to-CadQuery corpus~\citep{xie2025text}, which contains approximately \(110\text{k}\) CadQuery programs with geometric and semantic descriptions. Each program is treated as a reference design and analyzed using a deterministic CAD--CAE pipeline to obtain its maximum displacement \(\delta_{\max}\), \(99\)-th percentile von Mises stress \(\sigma_{\mathrm{P99}}\), and mass \(m\).

% Each valid reference is converted into an inverse-design task without exposing its exact geometry or simulation results. We redact exact numerical specifications while preserving high-level geometry semantics, and derive displacement, stress, and mass constraints from the corresponding reference measurements with a small feasibility margin. The resulting task provides the model with geometry semantics, coarse design bounds, physical conditions, and performance requirements, and asks it to generate a design satisfying the constraints while minimizing mass. Further details of the task construction and sanitization procedure are provided in the Appendix.

% After CAD--CAE annotation, \(77{,}067\) designs successfully pass the complete pipeline. We divide them into disjoint SFT and GRPO subsets using the same prompt format. The SFT subset contains \(11{,}024\) samples with sanitized reference CadQuery programs as supervision, while the GRPO subset contains \(66{,}043\) samples without target programs and is optimized solely through CAD--CAE verification.


\begin{table*}[t]
    \centering
    \small
    \caption{
    \textbf{Quantitative comparison on the full evaluation set of 3,373 tasks.}
    All results are reported as success rates (\%).
    Bold values indicate the best result among different settings of the same model.
    }
    \label{tab:model_comparison}
    \setlength{\tabcolsep}{5.0pt}
    \renewcommand{\arraystretch}{1.12}
    \resizebox{\textwidth}{!}{
    \begin{tabular}{llccccccc}
        \toprule
        \multirow{2}{*}{Model}
        & \multirow{2}{*}{Setting}
        & \multirow{2}{*}{Valid CAD$\uparrow$}
        & \multirow{2}{*}{BBox$\uparrow$}
        & \multirow{2}{*}{CAE$\uparrow$}
        & \multicolumn{3}{c}{Physical Constraints$\uparrow$}
        & \multirow{2}{*}{Joint$\uparrow$} \\
        \cmidrule(lr){6-8}
        & & & & 
        & Disp.
        & Stress
        & Physics
        & \\
        \midrule

        Claude Opus 5
        & Direct
        & 73.82 & 42.31 & 70.71
        & 46.52 & 49.87 & 8.42 & 7.29 \\

        GPT-5.6 Sol
        & Direct
        & 62.41 & 35.61 & 56.69
        & 46.58 & 35.75 & 16.87 & 15.51 \\

        DeepSeek-V4-Flash-Ascend
        & Direct
        & 63.33 & 50.52 & 60.45
        & 39.79 & 38.36 & 10.85 & 9.28 \\

        \midrule

        \multirow{3}{*}{Qwen3-4B}
        & Zero-shot
        & 18.56 & 11.41 & 17.55
        & 12.24 & 13.93 & 1.22 & 0.74 \\
        & SFT
        & 94.25 & 77.73 & 91.49
        & 63.09 & 63.27 & 39.19 & 36.64 \\
        & GRPO
        & \textbf{99.14} & \textbf{91.25} & \textbf{96.92}
        & \textbf{84.44} & \textbf{81.83}
        & \textbf{57.40} & \textbf{55.71} \\

        \midrule

        \multirow{3}{*}{Ministral-3-3B}
        & Zero-shot
        & 1.27 & 0.62 & 1.07
        & 0.68 & 0.65 & 0.24 & 0.12 \\
        & SFT
        & 95.55 & 78.33 & 92.68
        & 64.63 & 66.50 & 40.50 & 37.36 \\
        & GRPO
        & \textbf{98.99} & \textbf{94.22} & \textbf{97.45}
        & \textbf{84.32} & \textbf{81.20}
        & \textbf{59.29} & \textbf{58.05} \\

        \midrule

        \multirow{3}{*}{SmolLM3-3B}
        & Zero-shot
        & 6.91 & 2.91 & 6.55
        & 5.37 & 5.90 & 0.65 & 0.18 \\
        & SFT
        & 93.92 & 74.98 & 90.31
        & 58.85 & 60.72 & 31.28 & 28.73 \\
        & GRPO
        & \textbf{98.99} & \textbf{91.31} & \textbf{96.80}
        & \textbf{77.47} & \textbf{76.28}
        & \textbf{45.75} & \textbf{44.56} \\

        \midrule

        \multirow{3}{*}{StarCoder2-3B}
        & Zero-shot
        & 10.94 & 3.91 & 10.29
        & 8.81 & 8.98 & 1.04 & 0.56 \\
        & SFT
        & 93.74 & 74.50 & 90.16
        & 56.95 & 58.43 & 28.94 & 26.33 \\
        & GRPO
        & \textbf{98.96} & \textbf{92.05} & \textbf{97.15}
        & \textbf{76.67} & \textbf{74.12}
        & \textbf{43.14} & \textbf{41.65} \\

        \midrule

        \multirow{3}{*}{Qwen2.5-3B}
        & Zero-shot
        & 1.45
        & 0.47
        & 1.36
        & 1.01
        & 0.95
        & 0.21
        & 0.00 \\
        & SFT
        & 90.07 & 65.16 & 86.45
        & 55.68 & 57.75 & 25.32 & 21.76 \\
        & GRPO
        & \textbf{98.55} & \textbf{88.50} & \textbf{96.32}
        & \textbf{74.47} & \textbf{70.53}
        & \textbf{40.94} & \textbf{39.49} \\

        \bottomrule
    \end{tabular}
    }
    \vspace{-1em}
\end{table*}


\subsection{Implementation Details}
\label{sec:implementation}

From the \(66{,}043\) samples allocated to RL, we reserve \(3{,}373\) samples as a held-out evaluation set and use the remaining \(62{,}670\) prompts for GRPO training. The held-out samples are excluded from all policy updates and are used consistently for evaluating the SFT and GRPO checkpoints. All training and CAD--CAE evaluation experiments are conducted on a NVIDIA H800 GPU.

For GRPO, we use a training batch size of \(4\) prompts. For each prompt, the policy generates \(4\) independent completions, resulting in \(16\) rollout samples per batch. The four completions corresponding to the same prompt form a GRPO group and are evaluated independently by the deterministic CAD--CAE verifier described in Sec.~\ref{sec:verifier}. Their verification outcomes are then converted into the hierarchical advantages introduced above. We use a learning rate of \(5\times10^{-7}\) for policy optimization, with a sampling temperature of \(0.9\) and top-\(p\) of \(0.95\). For the hierarchical advantage, we instantiate ~\eqref{eq:hba} with
$
    \lambda_L = 1.0,\,
    \lambda_w = 0.25
$
and $\beta=0.01$ in~\eqref{eq:grpo}.


% \begin{figure}[t]
%     \centering
%     \includegraphics[width=\textwidth]{fig/result_new.pdf}
% \caption{
% \textbf{Qualitative comparison across training stages and model families.}
% Surface color indicates stress relative to the prescribed limit, while the bars
% summarize displacement, stress, and mass constraint satisfaction.
% Across model families, SFT improves CAD generation but can still violate geometric
% or physical constraints, whereas GRPO more consistently achieves joint feasibility.
% }
%     \label{fig:qualitative_comparison}
% \end{figure}



\subsection{Evaluation Metrics}
\label{sec:metrics}

We evaluate generated designs using seven metrics. \textbf{Valid CAD} measures whether the generated program successfully produces a valid CAD solid. \textbf{BBox} measures whether all three bounding-box dimensions satisfy the specified design bounds. \textbf{CAE} denotes successful finite-element analysis with valid physical measurements. \textbf{Displacement} (Disp.) and \textbf{Stress} report the pass rates for the prescribed maximum-displacement and stress constraints, respectively. \textbf{Physics} measures simultaneous satisfaction of the displacement, stress, and mass constraints. Finally, \textbf{Joint} measures whether both the geometric requirements and all physical constraints are satisfied.


\subsection{Comparison with the State of the Art}
\label{sec:sota}

We compare PhyCAD with three strong proprietary models under direct prompting,
including Claude Opus 5, GPT-5.6 Sol, and DeepSeek-V4-Flash-Ascend.
We further evaluate five open-source models under SFT and GRPO, with zero-shot
results reported where applicable, to examine the effectiveness of our training
pipeline across different model families.

Table~\ref{tab:model_comparison} and Fig.~\ref{fig:qualitative_comparison}
show a consistent progression across training stages.
SFT substantially improves CAD validity and CAE success, while GRPO further
improves geometric and physical constraint satisfaction across all evaluated
models. In particular, Joint success increases from
\(21.76\%\)--\(37.36\%\) after SFT to \(39.49\%\)--\(58.05\%\) after GRPO,
with corresponding gains in displacement and stress satisfaction.
The trained models also achieve substantially higher success on the complete
set of design requirements than proprietary models under direct prompting.
These results indicate that CAD--CAE verification provides effective supervision
beyond program validity and guides generation toward designs that better satisfy geometric and physical constraints.



\subsection{Ablation on hierarchical learning.}
Table~\ref{tab:training_stage_comparison} evaluates the effect of the hierarchical objective and staged optimization. Starting from the same SFT model, non-hierarchical GRPO improves CAD validity to $98.40\%$ but reduces joint feasibility from $21.76\%$ to $18.80\%$. In contrast, direct hierarchical GRPO achieves $98.55\%$ valid CAD and $39.49\%$ joint feasibility, nearly doubling the feasible designs relative to SFT. This indicates the gain comes from explicitly organizing feedback by feasibility level rather than from GRPO alone.

The staged results further show that optimizing geometric or physical objectives independently is insufficient: Geometry and Physics reach only $22.86\%$ and $27.63\%$ joint feasibility, respectively, improving over SFT by only $1.10$ and $5.87$ percentage points. Joint feasibility rises to $42.01\%$ only when the explicit Hierarchical objective is introduced. This is particularly important for physical feasibility, where satisfying the prescribed stress and displacement thresholds requires the model to prioritize constraint satisfaction before objective refinement. Although staged optimization yields a further $2.52$-point gain over direct hierarchical GRPO, it requires additional RL stages and substantially higher training cost.

% \subsection{Ablation on hierarchical learning.}
% Table~\ref{tab:training_stage_comparison} evaluates the effect of the hierarchical training objective and its staged optimization. Starting from the same SFT model, non-hierarchical GRPO substantially improves CAD validity to $98.40\%$, but does not improve constraint satisfaction, with joint feasibility decreasing from $21.76\%$ to $18.80\%$. In contrast, direct hierarchical GRPO jointly organizes geometric and physical feedback within a single reinforcement-learning objective, achieving $98.55\%$ valid CAD and $39.49\%$ joint feasibility. This comparison shows that the improvement primarily comes from the hierarchical formulation rather than from GRPO alone.
% The staged results further show that simply decomposing geometric and physical objectives into separate training stages is not sufficient to obtain the benefit of hierarchical learning. The Geometry and Physics objectives alone reach only $22.86\%$ and $27.63\%$ joint feasibility, respectively, substantially below the $39.49\%$ achieved by direct hierarchical GRPO. Joint feasibility increases to $42.01\%$ only after the explicit Hierarchical objective is introduced. Thus, sequentially optimizing individual constraint groups alone does not reproduce the benefit of the hierarchical formulation, which explicitly organizes geometric and physical feedback. While the staged strategy provides a further improvement over direct hierarchical GRPO, it requires additional RL stages and therefore incurs substantially higher training cost.










\begin{table}[t]
\centering
\caption{
\textbf{Ablation of hierarchical reinforcement learning.} Comparison of direct hierarchical GRPO, non-hierarchical GRPO, and staged training objectives.
}
\label{tab:training_stage_comparison}
\resizebox{\linewidth}{!}{
\begin{tabular}{llcccccc}
\toprule
\textbf{Training Strategy}
& \textbf{Stage}
& \textbf{Valid CAD} $\uparrow$
& \textbf{BBox} $\uparrow$
& \textbf{CAE} $\uparrow$
& \textbf{Physics} $\uparrow$
& \textbf{Joint} $\uparrow$
& \textbf{Code Tokens} $\downarrow$ \\
\midrule

\multirow{2}{*}{Baseline}
& Zero-shot
& 1.45 & 0.47 & 1.36 & 0.21 & 0.00 & 1002.41 \\

& SFT
& 90.07 & 65.16 & 86.45 & 25.32 & 21.76 & 201.06 \\

\midrule

Direct
& Non-hierarchical GRPO
& 98.40 & 60.90 & 83.60 & 23.40 & 18.80 & 199.78 \\

\midrule

Direct
& Hierarchical GRPO
& 98.55 & 88.50 & 96.32 & 40.94 & 39.49 & 195.46 \\

\midrule

\multirow{3}{*}{Training Objective}
& $\rightarrow$ Geometry GRPO
& 95.29 & 72.37 & 91.55 & 25.88 & 22.86 & 191.60 \\

& $\rightarrow$ Physics GRPO
& 94.72 & 75.27 & 91.40 & 30.33 & 27.63 & 188.03 \\

& $\rightarrow$ Hierarchical GRPO
& \textbf{99.23}
& \textbf{92.05}
& \textbf{97.45}
& \textbf{43.34}
& \textbf{42.01}
& \textbf{184.85} \\

\bottomrule
\end{tabular}
}
\vspace{-0.8em}
\end{table}


% \begin{table}[t]
% \centering
% \caption{
% \textbf{Comparison of direct and staged hierarchical reinforcement learning.}
% We compare direct hierarchical GRPO after SFT with a staged training strategy
% that progressively introduces geometry, physics, and hierarchical objectives.}
% \label{tab:training_stage_comparison}
% \resizebox{\linewidth}{!}{
% \begin{tabular}{llcccccc}
% \toprule
% \textbf{Training Strategy}
% & \textbf{Stage}
% & \textbf{Valid CAD} $\uparrow$
% & \textbf{BBox} $\uparrow$
% & \textbf{CAE} $\uparrow$
% & \textbf{Physics} $\uparrow$
% & \textbf{Joint} $\uparrow$
% & \textbf{Code Tokens} $\downarrow$ \\
% \midrule

% \multirow{2}{*}{Baseline}
% & Zero-shot
% & 1.45 & 0.47 & 1.36 & 0.21 & 0.00 & 1002.41 \\

% & SFT
% & 90.07 & 65.16 & 86.45 & 25.32 & 21.76 & 201.06 \\

% \midrule

% Direct
% & SFT $\rightarrow$ Hierarchical GRPO
% & 98.55 & 88.50 & 96.32 & 40.94 & 39.49 & 195.46 \\

% \midrule

% \multirow{3}{*}{Staged}
% & $\rightarrow$ Geometry
% & 95.29 & 72.37 & 91.55 & 25.88 & 22.86 & 191.60 \\

% & $\rightarrow$ Physics
% & 94.72 & 75.27 & 91.40 & 30.33 & 27.63 & 188.03 \\

% & $\rightarrow$ Hierarchical GRPO
% & \textbf{99.23}
% & \textbf{92.05}
% & \textbf{97.45}
% & \textbf{43.34}
% & \textbf{42.01}
% & \textbf{184.85} \\

% \bottomrule
% \end{tabular}
% }
% \vspace{-0.2em}
% \end{table}




\section{Discussion}
\paragraph{Conclusion.} Our work establishes a physics-aware Text-to-CAD setting that connects executable CAD generation with deterministic CAE verification. By combining a physics-aware dataset, a CAD--CAE verifier, and hierarchical reinforcement learning, the framework improves geometric and physical compliance across multiple models. These results demonstrate the potential of incorporating verifiable engineering feedback into CAD generation beyond geometric reconstruction.


\paragraph{Limitation and Future Work.} The current study has a few scope limitations. First, the prompts mainly encode physical constraints, while broader engineering intent is only weakly represented. Second, the dataset does not systematically cover material variations, though the same verification pipeline can be extended to different material settings. Future work could further introduce local hard dimensional constraints to enforce critical engineering requirements and reduce ambiguity in the generated geometry.


\bibliography{iclr2027_conference}
\bibliographystyle{iclr2027_conference}

\clearpage
\appendix
% \section{Implementations of Dataset Generation}
% \label{app:dataset}

% This appendix provides additional details of the dataset-generation
% procedure introduced in Sec.~\ref{sec:dataset}. We describe how the
% original Text-to-CadQuery samples are converted into physics-aware
% inverse-design tasks, including CAD--CAE annotation, automatic physical
% condition construction, design-domain generation, geometry semantic
% construction, and dataset partitioning. We also provide representative
% examples of the final model-visible prompts.

% \subsection{Reference CAD--CAE Annotation}
% \label{app:dataset_annotation}

% We start from the Text-to-CadQuery corpus~\citep{xie2025text}, which
% contains approximately \(110\text{k}\) CadQuery programs paired with
% geometric descriptions. Each source program is treated as a
% reference design: it provides a valid geometric realization from
% which an inverse-design task can be constructed, but its exact geometry
% is not exposed to the model.

% Each CadQuery program is first executed to construct the corresponding
% solid. Samples are discarded if the program cannot be executed, does
% not produce exactly one valid connected solid, or cannot be processed by
% the downstream geometry and simulation pipeline. Successfully
% constructed solids are exported and analyzed using the same deterministic
% finite-element pipeline.

% All samples are evaluated using three-dimensional small-strain linear
% elasticity in SI units. We use isotropic structural steel with Young's
% modulus \(2.1\times10^{11}\,\mathrm{Pa}\), Poisson's ratio \(0.3\), and
% density \(7850\,\mathrm{kg/m^3}\). For each successfully simulated
% reference design, we record the maximum nodal displacement
% \(\delta_{\max}\), the volume-weighted \(99\)-th percentile element
% von Mises stress \(\sigma_{\mathrm{P99}}\), and the total mass \(m\).
% These quantities are used only to construct the performance requirements
% of the inverse-design task and are not directly revealed to the model.


% \subsection{Automatic Physical Conditions}
% \label{app:physical_conditions}

% Because the source corpus contains objects with substantially different
% shapes and orientations, the fixture and loading regions are constructed
% automatically from the CAD geometry rather than manually specified for
% individual samples.

% For each solid, we first compute the centroid of every CAD surface. The
% pair of surface centroids with the largest Euclidean separation defines
% the principal analysis direction, which is normalized to obtain the
% functional analysis-axis unit vector. All surface centroids are then
% projected onto this axis. The projected extent of the solid is used to identify its two terminal
% surface groups. Surfaces whose centroid projections lie within \(5\%\)
% of the negative or positive end of the projected range are assigned to
% the corresponding terminal group. This surface-level grouping is used
% instead of selecting a single face, allowing geometries with multiple
% coplanar or fragmented terminal surfaces to be handled consistently.

% The negative terminal group is used as the fixture, with all
% translational degrees of freedom constrained,
% \(u_x=u_y=u_z=0\). The positive terminal group is used as the loading
% region. A uniform traction is distributed over this group such that its
% integrated resultant force has magnitude \(100\,\mathrm{kN}\) and is
% directed along the positive analysis axis.

% Both the analysis-axis vector and the resulting fixture and loading
% specifications are included in the task prompt. Therefore, although the
% boundary conditions are automatically derived from the reference
% geometry during dataset construction, they become explicit physical
% conditions that must also be satisfied by the generated design.


% \subsection{Inverse-Design Task Construction}
% \label{app:inverse_task}

% After a reference design has passed CAD construction and FEM analysis,
% we convert it into an inverse-design task. The conversion preserves the
% design family and physical requirements while deliberately removing the
% exact reference dimensions and simulation measurements.

% \paragraph{Coarse design domain.}
% We first compute the reference solid's global bounding-box extents
% \(e_x\), \(e_y\), and \(e_z\). Instead of exposing these values directly,
% each extent \(e\) is replaced by the enclosing interval formed by two
% adjacent powers of two:
% \begin{equation}
%     \mathcal{B}(e)
%     =
%     \left[
%         2^{\lfloor\log_2 e\rfloor},
%         2^{\lfloor\log_2 e\rfloor+1}
%     \right].
%     \label{eq:dataset_coarse_bound}
% \end{equation}
% The prompt therefore specifies only an admissible range for the global
% extent along each axis. For example, a reference extent lying between
% \(0.25\,\mathrm{m}\) and \(0.5\,\mathrm{m}\) is represented only by that
% interval rather than by its exact value.

% \paragraph{Performance requirements.}
% The displacement, stress, and mass requirements are derived from the
% corresponding reference measurements. For each measured quantity \(v\),
% we apply a \(10\%\) feasibility margin and round the resulting value
% upward to two significant digits:
% \begin{equation}
%     \tau(v)
%     =
%     \operatorname{RoundUp}_{2\mathrm{sig}}(1.10\,v).
%     \label{eq:dataset_threshold}
% \end{equation}
% The same procedure is applied independently to
% \(\delta_{\max}\), \(\sigma_{\mathrm{P99}}\), and \(m\).
% The resulting values become upper bounds on maximum displacement,
% \(99\)-th percentile von Mises stress, and total mass, respectively.

% The margin keeps the original reference design feasible under small
% numerical variations introduced by CAD processing, meshing, and FEM
% evaluation. Upward rounding additionally prevents the prompt from
% revealing the exact measurements of the reference design. 

% \subsection{Geometry Semantic Construction}
% \label{app:geometry_semantics}

% A central objective of the task construction is to preserve the intended
% geometric structure while preventing the model from directly recovering
% the original design parameters. We therefore construct the model-visible
% geometry semantics jointly from the source textual description and the
% CadQuery operation sequence.

% \paragraph{Numerical redaction.}
% We first sanitize the original textual description by removing numerical
% information that can reveal the reference geometry. This includes
% explicit dimensions, angles, ratios, percentages, quantities with units,
% and spelled-out numerical values. In contrast, non-numerical information
% describing the object category, shape composition, topology, feature
% types, feature counts, and spatial relationships is retained.

% For example, a source description may specify the exact diameter and
% height of a cylindrical feature. The sanitized description retains that
% a cylindrical feature exists and how it is positioned relative to the
% remaining geometry, but removes the original diameter and height. The
% generated program must determine such dimensions from the allowed design
% domain and physical requirements.

% \paragraph{Operation-derived semantics.}
% Textual descriptions alone do not always provide a sufficiently
% consistent representation of the underlying CAD construction. We
% therefore additionally inspect the CadQuery program and convert its
% modeling operations into dimension-free feature descriptions using a
% fixed operation-to-feature mapping. Numerical arguments of these
% operations are discarded.

% Table~\ref{tab:semantic_mapping} shows representative mappings used
% during this process. The mapping describes the geometric role of an
% operation rather than its parameters. For instance, a
% \texttt{cylinder} operation indicates the presence of a cylindrical
% feature, while a \texttt{cut} operation indicates a material-removal
% relationship. Sketches followed by extrusion retain the corresponding
% profile type and extrusion relation without preserving the original
% profile dimensions or extrusion distance.

% \begin{table}[t]
%     \centering
%     \small
%     \caption{
%     Representative CadQuery operation-to-feature mappings used to
%     construct dimension-free geometry semantics. Numerical parameters
%     associated with these operations are not exposed to the model.
%     }
%     \label{tab:semantic_mapping}
%     \setlength{\tabcolsep}{8pt}
%     \renewcommand{\arraystretch}{1.12}
%     \begin{tabular}{ll}
%         \toprule
%         \textbf{CadQuery construction} &
%         \textbf{Extracted geometry semantics} \\
%         \midrule
%         \texttt{box}
%         & rectangular-prismatic feature \\
%         \texttt{cylinder}
%         & cylindrical feature \\
%         polygonal sketch + \texttt{extrude}
%         & polygonal-prismatic feature \\
%         curved sketch + \texttt{extrude}
%         & curved-profile feature \\
%         \texttt{revolve}
%         & revolved feature \\
%         \texttt{cut}
%         & material-removal relation \\
%         \texttt{hole}
%         & hole or through-feature \\
%         \texttt{fillet}
%         & rounded-edge feature \\
%         \texttt{shell}
%         & hollow or shell construction \\
%         union / composition operations
%         & part-composition relation \\
%         \bottomrule
%     \end{tabular}
% \end{table}

% The sanitized source description and operation-derived information are
% then combined into a geometry-semantic description. The resulting
% description preserves the overall shape class, topology, feature types
% and counts, and relationships between individual components. Exact
% dimensions remain unspecified.

% For example, a reference composed of a hexagonal prism and a centered
% cylindrical top is described as requiring these two feature types, their
% relative placement, and their combination into a single solid, while the
% side lengths, cylinder diameter, and component heights remain hidden.
% Similarly, a design containing an extruded star-shaped profile combined
% with a cylindrical component preserves the corresponding feature
% composition without exposing the dimensions of either component.

% This construction makes the original reference a feasible realization
% of the requested design semantics, but not the uniquely specified
% solution. The model must select a concrete geometric realization that
% preserves the required structure while satisfying the design-domain and
% physical constraints.


% \subsection{Model-Visible and Private Information}
% \label{app:dataset_visibility}

% We maintain a strict separation between information used internally to
% construct each sample and information made available to the model. The
% model-visible prompt contains the dimension-free geometry semantics,
% coarse global bounding-box ranges, functional analysis axis, material
% properties, fixture and loading conditions, and the resulting
% displacement, stress, and mass limits.

% The exact reference dimensions are never included in the prompt.
% Likewise, the original CadQuery program, reference STEP model, FEM mesh,
% nodal displacement field, element stress field, and unrounded reference
% measurements remain private during generation. The reference CAD and FEM
% results therefore serve only as internal data for task construction and
% verification.

% This separation is particularly important for the inverse-design
% setting. The generated design is not evaluated by geometric similarity
% to the reference model. Instead, it must preserve the requested
% geometric semantics, remain inside the prescribed design domain, and
% satisfy the physical requirements through its own geometric realization.


% \subsection{Dataset Filtering and Partitioning}
% \label{app:dataset_split}

% After CAD execution, geometry processing, meshing, and FEM annotation,
% \(77{,}067\) source designs successfully pass the complete dataset
% construction pipeline.

% We partition these samples by sample ID into disjoint SFT and GRPO
% subsets. The SFT subset contains \(11{,}024\) samples. Each SFT sample
% pairs the constructed task prompt with its sanitized reference CadQuery
% program, allowing the model to first learn executable program
% construction and the required output format.

% The remaining \(66{,}043\) samples form the GRPO subset. Unlike SFT,
% these samples do not expose a target CadQuery program during
% reinforcement learning. Candidate programs are instead evaluated by the
% deterministic CAD--CAE verifier according to their executable validity,
% geometric compliance, physical feasibility, and design objective.

% Within the GRPO subset, \(62{,}670\) samples are used for reinforcement
% learning and \(3{,}373\) are reserved for evaluation. The SFT and GRPO
% sets contain no overlapping sample IDs. Importantly, the same task
% construction procedure and prompt format are used across the two stages;
% the difference lies only in whether a reference program is supplied as
% supervision or the generated program is evaluated through the verifier.


% \subsection{Representative Examples}
% \label{app:dataset_examples}

% Figure~\ref{fig:dataset_examples} presents representative samples
% generated by the complete pipeline. For illustration, each example
% shows the private reference geometry and its corresponding FEM von Mises
% stress field together with the final model-visible task prompt.

% The examples illustrate the distinction between the private reference
% design and the public inverse-design specification. High-level geometric
% structure is retained in the prompt, while exact dimensions are replaced
% by coarse admissible ranges. Similarly, the exact reference simulation
% results remain hidden and are converted into displacement, stress, and
% mass requirements. The model must therefore determine a new concrete
% set of dimensions through generation rather than reproduce the reference
% geometry directly.

\section{Implementation of Dataset Generation}
\label{app:dataset}

This appendix provides implementation details for the dataset-generation
procedure introduced in Sec.~\ref{sec:dataset}. We focus on the construction
of physical conditions, inverse-design constraints, and dimension-free
geometry semantics. 

\subsection{Reference CAD--CAE Annotation and Physical Conditions}
\label{app:dataset_annotation}

Each source CadQuery program is executed and retained only if it produces
exactly one valid connected solid that can be processed by the downstream
CAD--CAE pipeline. We analyze all retained designs using three-dimensional
small-strain linear elasticity in SI units, with isotropic structural steel
($E=2.1\times10^{11}$ Pa, $\nu=0.3$, and
$\rho=7850~\mathrm{kg/m^3}$). For each successfully simulated design, we
record the maximum nodal displacement $\delta_{\max}$, the volume-weighted
99-th percentile element von Mises stress $\sigma_{\mathrm{P99}}$, and the
total mass $m$.

Because the source geometries vary substantially in shape and orientation,
the fixture and loading regions are determined automatically. We compute the
centroid of every CAD surface and select the pair of surface centroids with
the largest Euclidean separation. The normalized direction between this pair
defines the functional analysis axis. All surface centroids are projected
onto this axis, and surfaces whose projections lie within $5\%$ of either end
of the projected range form the two terminal surface groups. This grouping
allows multiple coplanar or fragmented terminal surfaces to be handled
consistently.

The negative terminal group is fully fixed
($u_x=u_y=u_z=0$), while a uniform traction is applied to the positive
terminal group such that its integrated resultant force is $100$ kN along
the positive analysis axis. The resulting analysis axis, fixture definition,
and loading condition are explicitly included in the model prompt.

\subsection{Inverse-Design Task Construction}
\label{app:inverse_task}

Each successfully analyzed reference is converted into an inverse-design task
while withholding its exact dimensions and simulation measurements.

\paragraph{Coarse design domain.}
For each global bounding-box extent $e$, we expose only the enclosing interval
between adjacent powers of two,
\begin{equation}
    B(e)
    =
    \left[
    2^{\lfloor \log_2 e \rfloor},
    2^{\lfloor \log_2 e \rfloor+1}
    \right].
    \label{eq:bbox_domain}
\end{equation}
The model therefore receives a coarse admissible range along each axis rather
than the dimensions of the reference design.

\paragraph{Performance requirements.}
The displacement, stress, and mass limits are derived independently from
$\delta_{\max}$, $\sigma_{\mathrm{P99}}$, and $m$. Each reference measurement
is multiplied by $1.10$ and rounded upward to two significant digits to form
the corresponding upper bound. The margin keeps the reference design feasible
under small numerical variations from CAD processing, meshing, and FEM
evaluation, while rounding avoids exposing its exact simulation measurements.




\begin{figure*}[t]
    \centering
    \includegraphics[width=\textwidth]{fig/dataset_cad_fulu.pdf}
    \caption{
    \textbf{Representative physics-aware Text-to-CAD tasks.}
    Each example shows the private reference geometry and its FEM stress field
    together with the corresponding model-visible task. Exact reference
    dimensions and simulation measurements are withheld; the model receives
    dimension-free geometry semantics, coarse design bounds, physical
    conditions, and displacement, stress, and mass limits. It must generate
    an executable CadQuery design satisfying these requirements while
    minimizing mass.
    }
    \label{fig:dataset_gallery}
\end{figure*}




\subsection{Geometry Semantic Construction}
\label{app:geometry_semantics}

The geometry specification must preserve the intended design structure without
revealing the parameters of the reference CAD. We construct this specification
from both the source textual description and the CadQuery operation sequence.

\paragraph{Numerical redaction.}
We remove numerical information that may reveal the reference geometry,
including explicit dimensions, angles, quantities with
units, and spelled-out numerical values. We retain non-numerical information
describing the shape class, topology, feature types and counts, and spatial
or compositional relationships. Thus, for example, the existence and relative
placement of a cylindrical feature are preserved while its diameter and height
are withheld.

\paragraph{Operation-derived semantics.}
Because textual descriptions do not always capture the CAD construction
consistently, we additionally parse the CadQuery program and map its modeling
operations to dimension-free geometric features. Numerical operation arguments
are discarded. Representative mappings are shown in
Table~\ref{tab:operation_semantics}.

\begin{table}[t]
    \centering
    \caption{
    Representative CadQuery operation-to-feature mappings used to construct
    dimension-free geometry semantics. Numerical parameters are not exposed
    to the model.
    }
    \label{tab:operation_semantics}
    \small
    \begin{tabular}{ll}
        \toprule
        \textbf{CadQuery construction} & \textbf{Extracted geometry semantics} \\
        \midrule
        box                         & rectangular-prismatic feature \\
        cylinder                    & cylindrical feature \\
        polygonal sketch + extrude  & polygonal-prismatic feature \\
        curved sketch + extrude     & curved-profile feature \\
        revolve                     & revolved feature \\
        cut                         & material-removal relation \\
        hole                        & hole or through-feature \\
        fillet                      & rounded-edge feature \\
        shell                       & hollow or shell construction \\
        union / composition         & part-composition relation \\
        \bottomrule
    \end{tabular}
\end{table}

The sanitized textual description and operation-derived features are combined
into the final geometry-semantic specification. It preserves the overall shape
class, topology, feature composition, and relationships between components,
while leaving their exact dimensions unspecified. Consequently, the reference
CAD remains one feasible realization of the task rather than a uniquely
specified target geometry.

Representative tasks produced by the complete pipeline are shown in
Fig.~\ref{fig:dataset_gallery}, where the reference CAD and FEM results are
included only for illustration and are not provided to the model.

% \begin{figure*}[t]
%     \centering
%     \includegraphics[width=\textwidth]{fig/dataset_cad_fulu.pdf}
%     \caption{
%     \textbf{Representative physics-aware Text-to-CAD tasks.}
%     Each example shows the private reference geometry and its FEM von
%     Mises stress field together with the corresponding model-visible
%     task prompt. Exact reference dimensions and simulation measurements
%     are removed: the geometry is represented by dimension-free semantic
%     descriptions and coarse bounding-box ranges, while the reference
%     physical measurements are converted into displacement, stress, and
%     mass limits. The model must generate an executable CadQuery program
%     that preserves the specified geometric structure, satisfies all
%     physical requirements, and minimizes mass. The reference geometry
%     and FEM visualization are shown only for illustration and are not
%     provided to the model.
%     }
%     \label{fig:dataset_examples}
% \end{figure*}




\section{Hierarchical Balanced GRPO}
\label{app:hierarchical_grpo}

Algorithm~\ref{alg:hierarchical_grpo} summarizes the training procedure of our
Hierarchical Balanced GRPO. 

\begin{algorithm}[h]
\caption{\textbf{Hierarchical Balanced GRPO.}}
\label{alg:hierarchical_grpo}
\begin{algorithmic}[1]
\REQUIRE Policy $\pi_\theta$, reference policy $\pi_{\mathrm{ref}}$,
prompts $\mathcal{Q}$, group size $G$, weights $\lambda_L,\lambda_w,\beta$

\FOR{each training iteration}
    \STATE Sample a prompt batch $\mathcal{B}\subset\mathcal{Q}$

    \FOR{each prompt $q\in\mathcal{B}$}
        \STATE Sample $\{o_i\}_{i=1}^{G}\sim\pi_\theta(\cdot\mid q)$

        \STATE \textbf{Verification:}
        evaluate each $o_i$ with the CAD--CAE verifier to obtain
        level $L_i\in\{0,1,2,3\}$ and level-specific scores
        $\{s_{ik}\}_{k=1}^{K_{L_i}}$
        \STATE \hspace{\algorithmicindent}
        $L_i=0$: invalid CAD;\;
        $L_i=1$: geometry failure;\;
        $L_i=2$: physics failure;\;
        $L_i=3$: fully feasible

        \STATE \textbf{Between-level advantage:}
        $\bar L=\frac{1}{G}\sum_{j=1}^{G}L_j$,\quad
        $A_i^{\mathrm{level}}=L_i-\bar L$

        \STATE \textbf{Within-level PeerRank:}
        initialize $\widetilde A_i=0$ for all $i$
        \FOR{$\ell=0,1,2,3$}
            \STATE $I_\ell=\{i:L_i=\ell\}$
            \IF{$|I_\ell|>1$}
                \STATE Compute tie-aware centered ranks
                $r_{ik}$ of $\{s_{ik}\}_{i\in I_\ell}$
                for $k=1,\ldots,K_\ell$
                \STATE $\widetilde A_i
                =\sum_{k=1}^{K_\ell}w_{\ell k}r_{ik}$,
                \quad $\forall i\in I_\ell$
            \ENDIF
        \ENDFOR

        \STATE \textbf{Hierarchical advantage:}\quad
        $A_i=\lambda_L A_i^{\mathrm{level}}
        +\lambda_w\widetilde A_i$
    \ENDFOR

    \STATE Update $\theta$ using DAPO-style GRPO:
    $\mathcal{L}
    =\mathcal{L}_{\mathrm{DAPO}}(\theta;A)
    +\beta D_{\mathrm{KL}}(\pi_\theta\|\pi_{\mathrm{ref}})$
\ENDFOR

\RETURN Optimized policy $\pi_\theta$
\end{algorithmic}
\end{algorithm}


\end{document}
